% LEAD REPORT::ARC-0001::2026-10-07T14:57:38-07:00
\documentclass[11pt]{article}
\usepackage[margin=1in]{geometry}
\usepackage{amsmath,amssymb,amsthm,mathtools,hyperref}
\newtheorem{theorem}{Theorem}
\newtheorem{definition}[theorem]{Definition}
\newcommand{\SAT}{\mathrm{SAT}}
\newcommand{\R}{\mathcal R_{\SAT}}
\newcommand{\B}{\mathcal B}
\title{ARC-0001: Finite Existential Return Cycles}
\author{LEAD, ARC/WRAPPER research team}
\date{7 October 2026}
\begin{document}
\maketitle

\begin{abstract}
This report formalizes a bounded-horizon return-cycle version of the
achievable-region route. Individual improvement phases may expand time or
space; only their finite composite must contract a nonnegative envelope to a
fixed point strictly inside the Williams forbidden region. The result is an
exact conditional iteration lemma and does not establish the SAT-specific
return relation. It requires neither a universal implementation compiler nor
an input-length-growing recursion. P versus NP remains open.
\end{abstract}

\section{Status and claim card}
\begin{description}
\item[] \textbf{Claim ID:} ARC-RETURN-CYCLE1.
\item[] \textbf{Status before run:} CANDIDATE, described only informally in
\texttt{main4.tex}.
\item[] \textbf{Status after run:} PROVED/EXACT as a conditional
achievable-region lemma; SAT-specific finite return closure remains OPEN.
\item[] \textbf{Desired conclusion:} finite entry of one jointly witnessed
SAT time-space pair into
$\B=\{(t,s):t+s<\gamma\}$, where
$\gamma=2\cos(\pi/7)$.
\item[] \textbf{Exact quantifiers:} under $P=NP$, first
$\exists z_0\in\R$. The open premise asks for fixed $P\ge0,q\ge0$ such that
for every $z\in\R\setminus\B$ there exist a finite $m$ and points
$z=z^{(0)},\ldots,z^{(m)}=w$ in $\R$ with $w\le Pz+q$.
\item[] \textbf{Machine/model:} one fixed deterministic random-access model,
ordinary explicit SAT encoding, simultaneous worst-case time and space.
\item[] \textbf{SAME\_REALIZATION:} each exponent pair has both coordinates
witnessed by one uniform exact SAT machine. Coordinates may not be borrowed
from different machines.
\item[] \textbf{Weakest useful theorem:} the endpoint of a finite existential
return path obeys one stable affine envelope. No individual phase needs to
contract.
\item[] \textbf{Falsifiers:} a missing SAT successor theorem, a current-scale
call whose cost remains in the endpoint, a fixed point on or outside the
forbidden boundary, a negative matrix invalidating order propagation,
coordinate mixing, or input-length-dependent depth without a uniform
scheduling theorem.
\end{description}

\section{Corpus reviewed}
The accessible repository tree and directory listings were surveyed. The run
read root \texttt{main.tex}, \texttt{main2.tex}, \texttt{main3.tex},
\texttt{main4.tex}, \texttt{main5.tex}, \texttt{AGENTS.md},
\texttt{CODEX.md}, the complete repository \texttt{pnp-research} skill and its
synthesis/research-protocol references, the proof-spines protocol,
\texttt{spines.bib}, the newest accessible historical MAIN snapshot,
LEAD-0027, LEAD-0028, LEAD-0030, the RETURNMAP1/RETURNMAP2 records and finite
controls, the proof-spine tests/results, the task index, and the current corpus
reconciliation governance note. No \texttt{research/agents/proof},
\texttt{research/agents/research}, \texttt{research/agents/red}, or
\texttt{research/agents/write} directory existed at startup; the verified
initialization entry therefore remained the only predecessor record. A Google
Drive search for the combined ARC/WRAPPER manuscript names returned no
matching shared record, so the GitHub blobs were treated as authoritative.

\section{Duplication and contradiction audit}
The original ARC manuscript already noted that a multi-stage cycle could be
used, but did not state its quantifiers, composite affine parameters, or proof.
RETURNMAP1/2 concern a source-specific ordinary-CNF return and remain open at
that construction layer. They do not duplicate this existence-level lemma.
The root synthesis already distinguishes fixed finite recursion from a
growing-depth schedule and forbids borrowing resource coordinates. No
contradiction with a verified project result was found.

Relative to root \texttt{main.tex}, this run changes no executive claim. It
formalizes a conditional route only in \texttt{main4.tex}. The SAT-specific
inner transformation, model-matched endpoint application, nonrelativizing
hinge, and final separation remain open. No prior component was reactivated or
promoted beyond its verified scope.

\section{Primary literature and transfer audit}
Williams's STACS 2010 Theorem 3.9 and surrounding deterministic
time-space statement give the model-matched endpoint
$t(n)s(n)=\Omega(n^{2\cos(\pi/7)-o(1)})$ for SAT on the declared RAM model,
with big-$\Omega$ used in the infinitely-often sense \cite{williams2010}.
This theorem supplies the forbidden endpoint; it does not produce an
achievable-region successor.

Williams's STOC 2024 Theorem 1.1 derives a faster small-circuit
Circuit-SAT algorithm from a large-circuit algorithm, and Theorem 1.4 proves an
existential-to-universal equivalence for the number of inputs of Large Circuit
SAT. Its proof explicitly uses finitely many stages for every fixed parameter.
Section 5 asks whether self-improvement can extend to the $P=NP$ regime and
states that current arguments encounter bottlenecks \cite{williams2024}.
This is a directly relevant precedent for staged algorithm existence, not a
transfer of the missing ordinary-SAT closure.

The control-theory check considered common polynomial Lyapunov functions for
uniform asymptotic stability under arbitrary switching
\cite{mason2004}. The source proves existence of such functions and the lack
of a uniform degree bound across all stable systems. Classification:
ANALOGY/CANDIDATE only. Arbitrary switching quantifies over every switching
signal, while ARC selects one existential successor among uniform SAT
algorithms; the physical trajectory and SAT-achievable-region objects do not
match. No control theorem was imported as a SAT theorem.

\section{Finite return-cycle theorem}
\begin{definition}
Let $P$ be entrywise nonnegative and $q\ge0$. A finite $(P,q)$-return
witness for $z\in\R$ is a finite chain
\[
z=z^{(0)},z^{(1)},\ldots,z^{(m)}=w,
\qquad z^{(i)}\in\R,
\qquad w\le Pz+q.
\]
Intermediate points may expand in either coordinate.
\end{definition}

\begin{theorem}[ARC-RETURN-CYCLE1]
Assume $\R$ contains a finite point. Suppose there are fixed $P\ge0,q\ge0$
such that $\rho(P)<1$, $z_*=(I-P)^{-1}q\in\B$, and every
$z\in\R\setminus\B$ has a finite $(P,q)$-return witness. Then
$\R\cap\B\ne\varnothing$.
\end{theorem}
\begin{proof}
Choose finite $z_0\in\R$. While the current endpoint is outside $\B$, choose
the endpoint $z_{r+1}$ of one return witness. Let $u_0=z_0$ and
$u_{r+1}=Pu_r+q$. Because $P$ is nonnegative, induction gives
$z_r\le u_r$. Also
\[
u_r=z_*+P^r(z_0-z_*)\longrightarrow z_*.
\]
The forbidden set $\B$ is open and downward closed in the nonnegative
quadrant. Hence $u_r\in\B$ for all sufficiently large finite $r$, and then
$z_r\in\B$. The proof selected finitely many return endpoints. Every selected
return chain is finite, so their concatenation is finite. No input-length
dependent family is created.
\end{proof}

For fixed phases
\[
z^{(i+1)}\le M_i z^{(i)}+d_i
\quad (i=0,\ldots,m-1),
\]
with nonnegative $M_i,d_i$, order propagation gives composite parameters
\[
P=M_{m-1}\cdots M_0,
\qquad
q=d_{m-1}+\sum_{j=0}^{m-2}M_{m-1}\cdots M_{j+1}d_j.
\]
It is therefore sufficient to prove stability and the strict endpoint margin
for the composite, even if each phase expands one resource.

\section{Bounded model and falsification}
The reproducible command was
\[
\texttt{python experiments/proof\_spines/arc\_return\_cycle\_test.py}.
\]
The program used exact rational arithmetic for the affine maps, seed
20261007, a 25-second cooperative budget, and finished 277 checks in about
0.007 seconds. The test used
\[
M_0=\operatorname{diag}(2,1/5),\qquad
M_1=\operatorname{diag}(1/5,2),\qquad
d_0=d_1=(1/20,1/20)^T.
\]
Each phase matrix has spectral radius two, but their composite is
$P=(2/5)I$ with $q=(3/50,3/20)^T$ and fixed point
$z_*=(1/10,1/4)^T$. All 65 seeded starting points entered the forbidden
half-plane in at most four cycles; this is a finite arithmetic check only.

\begin{description}
\item[] \textbf{Transient expansion:} the first phase expanded $t+s$ on the
recorded trace, while the complete cycle contracted toward $z_*$.
\item[] \textbf{Strict-decrease control:} the exact sequence
$\gamma+2^{-r}$ stayed above the boundary through all tested indices and
symbolically for every finite $r$.
\item[] \textbf{SAME\_REALIZATION control:} taking the time coordinate from
$(1.2,10)$ and the space coordinate from $(10,0.1)$ creates the fictitious
pair $(1.2,0.1)$ inside $\B$; neither machine witnesses that pair.
\item[] \textbf{Padding control:} padding can make an exponent $12$ appear as
$6,3,$ or $1.5$ under alternate length measures. This does not give an
algorithm on the original explicit SAT encoding.
\item[] \textbf{Order control:} a non-downward target set invalidates the
inference from an envelope point to a dominated achieved point.
\item[] \textbf{Growing-depth control:} for
$\lambda_N=1-1/\log_2N$, the tested depth needed to move from 10 below
$\gamma$ grew from 27 at $\log_2N=16$ to 438 at $\log_2N=256$. The finite
existential theorem does not cover such $J(N)\to\infty$ schedules.
\item[] \textbf{Easy-family control:} no SAT transformation was instantiated.
Horn, dual-Horn, 2-SAT, affine/Tseitin, bounded-treewidth, and independent
product families therefore cannot promote the abstract lemma; easy closure on
those subclasses would not imply closure for general SAT.
\end{description}

The code is
\path{experiments/proof_spines/arc_return_cycle_test.py}; raw output is
\path{results/raw/proof_spines/2026-10-07_145738_arc_return_cycle.json}.

\section{Barrier audit and exact open arrow}
ARC-RETURN-CYCLE1 is an order/convergence lemma. It relativizes whenever its
premise is assumed for the relativized achievable region, and algebraic-oracle
replacement does not disturb the affine calculation. Consequently it cannot
be the required nonrelativizing or nonalgebrizing hinge. DEAD, COPY, DELAY,
INTERP, RECOMP, source rereading, and semantic-oracle access remain attacks on
any proposed SAT-specific phase map, not on the conditional iteration proof.

The highest-value open arrow is now the weakest SAT-specific return statement:
\[
\boxed{
\forall z\in\R\setminus\B\;\exists\text{ a finite chain of uniform SAT
machines whose endpoint obeys one fixed stable composite envelope}.}
\]
This permits transient phase expansion and different programs at different
stages. It still requires exact general-SAT correctness, joint resource
accounting, strict endpoint margin, recursive eligibility, and model matching.
No such SAT-specific chain is proved.

\section{Manuscript changes and peer handoff}
Root \texttt{main4.tex} now contains ARC-RETURN-CYCLE1, its phased affine
corollary, the exact quantifier distinction, the two-phase example, and the
test record. Root \texttt{main5.tex} received a dated no-status-change review:
W1--W6 remain open. \texttt{spines.bib} now contains verified entries for the
2024 self-improvement paper and the switched-systems analogy. Root
\texttt{main.tex} was not edited because this team owns the two route
manuscripts and the new result does not change the executive proof status.

\begin{description}
\item[] \textbf{Handoff to PROOF:} search for one exact two-phase WRAPPER
candidate where phase one may trade time for space and phase two may trade
space for time. Report the actual composite $P,q$, every current-scale call,
and whether the endpoint is a uniform exact SAT machine on the original
encoding. Do not require either phase to contract separately, and do not use a
growing-depth schedule unless its effective errors and access costs are proved.
\end{description}

\bibliographystyle{plain}
\bibliography{spines}
\end{document}
% END_OF_VERIFIED_REPORT::LEAD::ARC-0001